\section{Prufer 序列}

一棵 $n$ 个点的带标号无根树与一个长度为 $n-2$ 的序列一一对应，该序列即 Prufer 序列。常用于无根树的计数、随机生成树等。编码与解码均为 $\mathcal{O}(n)$。

\begin{lstlisting}
struct Prufer{
    static std::vector<int> encode(std::vector<std::vector<int>>& G,int n){
        std::vector<int> fa = transfromFa(G,n);
        return encode(fa,n);
    }
    static std::vector<int> encode(std::vector<int>& fa,int n){
        std::vector<int> deg(n+1);
        for (int i = 1;i < n;i++) deg[fa[i]]++;
        std::vector<int> res(n-2);
        for (int i = 1, ptr = 1;i <= n-2;i++){
            while (deg[ptr] != 0) ptr++;
            res[i-1] = fa[ptr];
            deg[ptr]--;
            deg[fa[ptr]]--;
            if (deg[fa[ptr]] == 0 && fa[ptr] < ptr) ptr = fa[ptr];
        }
        return res;
    }
    static std::vector<int> decode(std::vector<int> a,int n){
        std::vector<int> deg(n+1),fa(n+1);
        for (int i = 0;i < a.size();i++) deg[a[i]]++;
        a.push_back(n);
        for (int cur = 1,i = 0,tmp = -1;i < a.size();i++){
            int u;
            if (tmp != -1) {
                u = tmp;
                tmp = -1;
            } else {
                while (deg[cur] > 0) cur++;
                u = cur;
                cur++;
            }
            fa[u] = a[i];
            deg[a[i]]--;
            if (deg[a[i]] == 0 && cur > a[i]) tmp = a[i];
        }
        return fa;
    }
private:
    static std::vector<int> transfromFa(std::vector<std::vector<int>>& G,int n){
        std::vector<int> fa(n+1);
        auto dfs = [&](this auto&& dfs,int u,int pa)->void {
            fa[u] = pa;
            for (int v : G[u]) {
                if (v == pa) continue;
                dfs(v,u);
            }
        };
        dfs(n,n);
        return fa;
    }
};
\end{lstlisting}

